% page 601 du fac-simile (image p-600 du scan)
% bloc 162.6 x 259.3 mm ; marge gauche 39.4 mm
\pgc{17.780mm}{0.000mm}{
\l{}
\l{}
\l{}
\l{}
\l{\cel{53}593}
\l{}
\l{\sou{treso}. Kordono plata, ek fili (lana, silka, e c.) od ek palio}
\l{\cel{3}o hari interplektita. - DEFIRS.}
\l{}
\l{\sou{tresto}. Planko sur quar pedi, uzata por subtenar longa tablo,}
\l{\cel{3}eshafodo,\cel{2}+stej (ceneyo) feriala. - EF.}
\l{}
\l{\sou{trezoro}. I. Asemblajo ek moneto-peci ora, arjenta, ek objekti}
\l{\cel{3}tre precoza - rezervita, celita od enterigita. II. (\sou{yuro-}}
\l{\cel{3}\sou{cienco}).Omna kozo celita ed enterigita quan ulu deskovras}
\l{\cel{3}hazarde, e pri qua nulu povas justifikar proprieto-yuro.}
\l{\cel{3}III. Ensemble de la revenui di stato. - IV. Kolektajo ek re-}
\l{\cel{3}liquii, vazi, ornivi tre precoza, qui Konservesas en kirko.}
\l{\cel{2}- V. (\sou{metaf.})\cel{2}Irgo quo egardesas kom tre precoza. (\sou{anke pri}}
\l{\cel{3}\sou{persono}). - DEFIS.}
\l{}
\l{\sou{tri}. (\sou{aritm.})\cel{2}La cifro \textquotedbl{}3\textquotedbl{}. - DEFIRS.}
\l{}
\l{\sou{tri.}\cel{2}Prefixo ciencala : Qua havas tri...}
\l{}
\l{\sou{triangulo}. I. (\sou{geom.}) Figuro geometriala kun tri anguli - parto}
\l{\cel{3}de planajo qua havas kom limiti tri rekti qui intersekas}
\l{\cel{3}duope. - II. Objekto qua havas la sama formo kam triangulo:}
\l{\cel{3}(a) (\sou{religio kristana}).Simbolo di la Triuno, (b) Emblemo di}
\l{\cel{3}la framasonaro ; (c) (\sou{muziko})\cel{2}Instrumento ek stalo, quan onu}
\l{\cel{3}resonigas per vergeto stala. - DEFIRS.}
\l{}
\l{\sou{triaso}. (\sou{geol.}) Sulo-strato ek marni, kalkajo, greso, qua su-}
\l{\cel{3}cedas liaso segun la docens-ordino. -}
\l{}
\l{\sou{tribono}. (en\cel{1}\sou{la Roma antiqua)}. 1.\cel{1}\sou{Tribono di la soldati}\cel{1}:\cel{2}Chefo}
\l{\cel{3}qua komandis legiono, alterne kun kin altri. - II. Tribono}
\l{\cel{3}di\cel{1}\sou{la plebeyi}\cel{1}:\cel{2}Judiciisto qua defensis la interesti di la}
\l{\cel{3}plebeyani. - DEFIRS.}
\l{}
\l{\sou{tribuo}. Grupo de familii qua formacas socio neperfekte orga-}
\l{\cel{3}nizita. - DEFIS.}
\l{}
\l{\sou{tribular}. (\sou{trans., per)}.\cel{2}Tormentar (lu) psikale, mentale.}
\l{\cel{3}- DEFIS.}
\l{}
\l{\sou{tribuno}. I. Loko situita kelke-alte, quaze estrado de ube la}
\l{\cel{3}oratoro diskursis a la plebeyani sur placo (en la epoki}
\l{\cel{3}antiqua) - de ube la oratoro diskursas en asemblitaro}
\l{\cel{3}deliberanta (nia-epoke). - II. Galerio ube onu rezervas}
\l{\cel{3}plasi en kirko, en asemblitaro deliberanta. - DEFIRS.}
\l{}
\l{\sou{tribunalo}.\cel{2}(\sou{yurocienco)}. Asemblitaro de judiciisti qui}
\l{\cel{3}+seancas (kunsidas) ensemble. - DEFIRS.}
\l{}
\l{\sou{tributar}. (\sou{netrans., ad)}.\cel{2}(\sou{aludante populo vinkita)}. Pagar}
\l{\cel{3}la pekunio-quanto, livrar la vari o precozaji, quin la}
\l{\cel{3}populo vinkinta impozis a sua dominacati. - DEFIS.}
\l{}
\l{\sou{tricepso}. (\sou{anat.)}\cel{3}Muskulo ek tri faski. - DEFS.}
}
